\chapter{Fracture Management}
\index{Fracture Management}

\section*{Definition and Overview}
Fracture management is the process of treating a broken bone, from the emergency care at the scene through to the restoration of full function. The goals are to realign the bone if displaced, hold it still while it heals, relieve pain, prevent complications, and rehabilitate the surrounding muscles and joints. Management ranges from simple measures for minor fractures, such as a splint or buddy taping, to surgery with plates and screws for displaced or unstable fractures. The choice of treatment depends on the bone, the pattern of the break, the patient's age and health, and the demands of their life and work. This chapter explains how fractures are managed and what recovery involves.

\section*{Epidemiology}
Fractures are among the most common injuries presenting to emergency departments, affecting around 1 in 15 people each year. The wrist, ankle, hip, and arm are the most commonly fractured sites, and the causes differ with age: sport and road trauma in young people, and falls in older adults, whose bones are weakened by osteoporosis. Hip fractures in older adults are particularly serious, with substantial mortality and loss of independence. Most fractures heal well with appropriate treatment, and the aim of modern management is a quick return to normal function.

\section*{Aetiology and Pathophysiology}
\begin{itemize}
    \item \textbf{Causes:} falls, direct blows, twisting injuries, road trauma, sport, and, in weakened bone, minor forces
    \item \textbf{Types:} simple (two pieces), comminuted (many pieces), displaced (ends out of alignment), open (through the skin), and greenstick (incomplete, in children)
    \item \textbf{Healing:} a fracture triggers an inflammatory response, then new bone (callus) forms over weeks and remodels over months
    \item \textbf{Stability:} fractures that are well aligned and stable heal without surgery; unstable or displaced fractures need fixation to hold the alignment
    \item \textbf{Blood supply:} healing depends on blood flow; bones with poor supply, such as the scaphoid and femoral head, heal slowly
    \item \textbf{The role of immobilisation:} a cast, splint, or internal fixation holds the bone still while callus forms
    \item \textbf{Soft tissue injury:} muscles, ligaments, and skin are often injured with the bone and need their own care
    \item \textbf{Children versus adults:} children's bones heal faster and remodel better; the growth plate is a special concern
\end{itemize}

\section*{Clinical Features}
\begin{itemize}
    \item Pain at the injury site, made worse by movement or weight bearing
    \item Swelling, bruising, and tenderness
    \item Deformity, shortening, or abnormal movement of the limb
    \item Inability to use the limb normally
    \item In open fractures: a wound over the fracture site
    \item Numbness, tingling, or a pale, cold limb, which indicates nerve or blood vessel injury
    \item In children: refusal to use the limb, and tenderness over a bone
    \item After treatment: pain settling, and gradual return of movement
\end{itemize}

\section*{Differential Diagnosis}
Injury symptoms have several possible causes.
\begin{itemize}
    \item Sprains, strains, and ligament injuries
    \item Dislocations, in which bones are out of joint
    \item Contusions and muscle injuries
    \item Tendon injuries and ruptures
    \item Stress fractures, which cause gradual pain without a specific injury
    \item In older adults, the possibility of pathological fracture through a tumour or weakened bone
    \item Imaging confirms the diagnosis
\end{itemize}

\section*{When to Seek Medical Care}
Any injury with severe pain, deformity, swelling, or inability to bear weight or use the limb should be assessed, with X-rays where indicated. Minor injuries that are still painful or swollen after a few days also need review.

\textbf{Contact your local emergency services for:}
\begin{itemize}
    \item Open fractures, with the bone through the skin
    \item A pale, cold, numb, or blue limb beyond the injury
    \item Severe deformity or an obviously dislocated joint
    \item Falls from height with back or hip pain
    \item Head injury with the fracture
    \item Inability to move or loss of sensation beyond the fracture
\end{itemize}

\section*{Diagnosis and Investigations}
\begin{itemize}
    \item \textbf{History and examination:} the mechanism of injury and a careful assessment of the limb, including pulses and sensation
    \item \textbf{X-rays:} two views are standard, and the joint above and below the fracture is included
    \item \textbf{CT scan:} for complex fractures, especially of the pelvis, heel, and joints
    \item \textbf{MRI:} for suspected stress fractures, ligament injury, or bone viability
    \item \textbf{Assessment of the neurovascular status:} essential for displaced and open fractures
    \item \textbf{Fracture clinic review:} follow-up imaging to confirm the bone is healing in good alignment
\end{itemize}

\section*{Management}
\begin{itemize}
    \item \textbf{First aid:} immobilise the limb, apply ice, elevate, and seek medical care; do not attempt to straighten a deformed limb
    \item \textbf{Reduction:} displaced fractures are realigned, under sedation or anaesthesia, before immobilisation
    \item \textbf{Non-operative treatment:} casts, splints, and braces hold the bone while it heals; this is used for stable, well-aligned fractures
    \item \textbf{Operative treatment:} plates, screws, rods, or external fixators hold displaced, unstable, or open fractures; surgery allows earlier movement
    \item \textbf{Pain relief:} simple analgesics, and stronger medicines in the acute phase
    \item \textbf{Weight-bearing restrictions:} many leg fractures require a period without full weight bearing
    \item \textbf{Physiotherapy:} essential after the bone heals to restore movement, strength, and function
    \item \textbf{Treat the whole patient:} manage osteoporosis, smoking, and nutrition, which affect healing
    \item \textbf{Complication watch:} plaster care, wound checks after surgery, and monitoring for signs of infection or poor healing
\end{itemize}

\section*{Complications}
\begin{itemize}
    \item Delayed union and non-union, in which the bone fails to heal
    \item Malunion, healing in a poor position
    \item Infection, especially after open fractures and surgery
    \item Nerve and blood vessel injury
    \item Compartment syndrome, a painful emergency caused by pressure within a limb
    \item Deep vein thrombosis after leg fractures and immobility
    \item Stiffness, muscle wasting, and chronic pain
    \item Complications of anaesthesia and surgery
\end{itemize}

\section*{Prognosis and Outcomes}
The vast majority of fractures heal completely, and most people return to their previous activities within weeks to months. Simple fractures in children heal quickly and remodel well. Surgical fixation allows earlier movement and better outcomes for displaced fractures, and rehabilitation is the key to regaining full function. Hip fractures in older adults carry the most serious outlook, but prompt surgery and rehabilitation improve recovery. With good treatment and rehabilitation, most fractures leave no lasting disability.

\section*{Prevention and Self-Care}
\begin{itemize}
    \item Maintain strong bones with calcium, vitamin D, and weight-bearing exercise
    \item Prevent falls at home: remove hazards, improve lighting, and use handrails
    \item Use protective equipment in sport, and seatbelts in cars
    \item Build up exercise loads gradually to prevent stress fractures
    \item Do not smoke, as it delays bone healing
    \item Follow the plaster and weight-bearing instructions after treatment
    \item Attend fracture clinic reviews and physiotherapy
    \item Seek prompt review for pain, swelling, or a change in the limb during recovery
\end{itemize}

\section*{Sources and Further Reading}
The following reputable sources provide further information for healthcare students and patients seeking reliable, current advice about fracture management.
\begin{itemize}
    \item Healthdirect Australia. \textit{Fractures.} https://www.healthdirect.gov.au/fractures
    \item Australian Orthopaedic Association. \textit{Fracture care.} https://www.aoa.org.au/
    \item Better Health Channel. \textit{Fractures.} https://www.betterhealth.vic.gov.au/
    \item NHS. \textit{Overview of fractures.} https://www.nhs.uk/conditions/fractures/
    \item Mayo Clinic. \textit{Fractures.} https://www.mayoclinic.org/
    \item World Health Organization. \textit{Falls prevention.} https://www.who.int/
\end{itemize}
